\documentclass{beamer}
\usetheme{Madrid}


\usepackage{comment}
\usepackage{amsmath, amssymb}  
\usepackage{amsthm}  
\usepackage{mathtools}  
\usepackage{complexity}

\newcommand{\cetal}{\textit{et al.\@}}  % shortcut for etal to be used when no space is desired
\newcommand{\etal}{\textit{et al.\@\ }}  % shortcut for etal to be used in regular text
\newcommand{\nrank}{\operatorname{rank}_+}
\newcommand{\rank}{\operatorname{rank}}
\newcommand{\diag}{\operatorname{diag}}
\newcommand{\bits}{\{0,1\}}
\newcommand{\abs}[1]{\left|#1\right|}
\newcommand{\conv}{\operatorname{conv}}
\newcommand{\suppmat}{\operatorname{suppmat}}
\newcommand{\supp}{\operatorname{supp}}
\newcommand{\ol}[1]{\overline{#1}}
\newcommand{\ul}[1]{\underline{#1}}
\newcommand{\oul}[1]{\overline{\underline{#1}}}
\newcommand{\xc}{\operatorname{xc}}
\newcommand{\TSP}{\operatorname{TSP}}
\newcommand{\STAB}{\operatorname{STAB}}
\newcommand{\CUT}{\operatorname{CUT}}
\newcommand{\COR}{\operatorname{COR}}
\renewcommand{\R}{\mathbb{R}}

\author[Livanos, Torres]{Vasilis Livanos and Manuel Torres}
\title[Exp. Lower Bounds for Polytopes]{Exponential Lower Bounds for Polytopes in Combinatorial Optimization}
\institute[UIUC]{University of Illinois at Urbana-Champaign}
\date{May 2, 2018}

\begin{document}

\begin{frame}
\titlepage
\end{frame}


\begin{frame}
\frametitle{Main Question}

\pause
\begin{itemize}
\item Let's prove $\P = \NP$!
\pause
\item Idea: Can we write a polynomial-size LP (Linear Program) for an \NP-Complete problem (e.g. $\lang{TSP}$)?
\pause
\item Another idea: Encode a problem as a polytope that is the convex hull of the vertices corresponding to the solutions of the problem.
\end{itemize}
\end{frame}


\begin{frame}
\frametitle{Necessary Background}

\pause
\begin{definition}[Extended Formulation]
Let $P = \{x : Ax \leq b\}$ and $Q = \{(x,y) : Bx + Cy \leq d\}$. Then $Q$ is an \emph{extended formulation (EF)} of $P$ if and only if
\[
P = \{x : \exists y , (x,y) \in Q\}
\]
The size of an extended formulation is the number of facets (faces of maximal dimension) of $Q$.
\end{definition}

\pause
\begin{definition}[Extension Complexity]
The \emph{extension complexity} of $P$, denoted $\xc(P)$, is the minimum size EF of $P$.
\end{definition}

\end{frame}


\begin{frame}
\frametitle{Extended Formulations Can Help}

\pause
\begin{itemize}
\item There are polytopes that have exponential size (i.e. number of facets), but have an EF with polynomial size.
\pause
\begin{enumerate}
\item The permutahedron polytope, which is the convex hull of the permutation group on $\{1, 2, \cdots, n\}$.
\pause
\item The spanning tree polytope.
\pause
\item The matching polytope for planar graphs.
\end{enumerate}
\pause
\item Thus, even though these polytopes have exponential size, we can write a polynomial-size LP for them, through their EF.
\end{itemize}
\end{frame}


\begin{frame}
\frametitle{Prior Work}

\pause
\begin{theorem}[Yannakakis, '91]
Every symmetric LP for TSP has extension complexity $2^{\Omega(n)}$. Thus, it also has an exponential number of inequalities.
\end{theorem}

\pause
\begin{corollary}
No symmetric LP can be used to solve TSP in polynomial time.
\end{corollary}

\pause
\begin{itemize}
\item Can we use an asymmetric LP to solve TSP in polynomial time?
\end{itemize}

\end{frame}


\begin{frame}
\frametitle{More Technical Background}

\pause
\begin{definition}[Nonnegative Rank]
Let $M \in \R_+^{m \times n}$. The \emph{nonnegative rank} of $M$ is defined as $\nrank(M) = \min\{r : M = TU, T \in \R_+^{m \times r}, U \in \R_+^{r \times n}\}$.
\end{definition}

\pause
\begin{definition}[Slack Matrix]
Let $A \in \R^{m \times d}$, $b \in \R^m$, and $V = \{v_1, \ldots, v_n\} \subseteq \R^d$. Let $P = \{x \in \R^d : Ax \le b\} = \conv(V)$. Then $S \in \R_+^{m\times n}$, where each entry is defined as $S_{ij} = b_i - A_i v_j$ with $i \in [m]$, $j \in [n]$, is the \emph{slack matrix} of $P$ with respect to $Ax \le b$ and $V$.
\end{definition}

\end{frame}


\begin{frame}
\frametitle{Yannakakis's Factorization Theorem}

\pause
\begin{theorem}[Factorization Theorem, Yannakakis '91]
Let $P = \left\{ x \in \R^d \mid Ax \leq b \right\} = \conv(V)$ be a polytope, and $S$ be its slack matrix with respect to $Ax \leq b$ and $V$. Then
\[
\xc(P) = \nrank(S)
\]
\end{theorem}

\pause
\begin{itemize}
\item If we want to lower bound $\xc(P)$, it suffices to lower bound $\nrank(S)$.
\end{itemize}

\end{frame}


\begin{frame}
\frametitle{Lower Bound on Nonnegative Rank}

\pause
\begin{definition}[Support Matrix]
The {\em support matrix} of a matrix $M$ with nonnegative entries is
\[
\suppmat(M)_{ij} = \begin{cases}
0 & M_{ij} = 0 \\
1 & M_{ij} \neq 0
\end{cases}
\]
\end{definition}

\pause
\begin{theorem}[Yannakakis, '91]
Let $M$ be a matrix with nonnegative entries, and $\suppmat(M)$ its support matrix.
Then, $\nrank(M)$ is lower bounded by the rectangle covering bound for $\suppmat(M)$.
\end{theorem}

\pause
\begin{itemize}
\item This provides a connection between $\xc(P)$ and communication complexity.
\end{itemize}

\end{frame}


\begin{frame}
\frametitle{Reminder: Rectangle Covers}

\pause
\begin{definition}[Rectangle Cover]
Let $M \in \bits^{2^n \times 2^n}$ with indices corresponding to bit strings. A \emph{rectangle} is a subset of $\bits^n \times \bits^n$. We say that $R \subseteq \bits^n \times \bits^n$ is a \emph{$b$-monochromatic rectangle} for $f$ if $M_{xy} = b$ for all $(x,y) \in R$. We say that a set $\mathcal{R}$ of $b$-monochromatic rectangles is a \emph{$b$-rectangle cover} if $\{(x,y) \in \bits^n \times \bits^n : M_{xy} = b\} \subseteq \bigcup_{R \in \mathcal{R}} R$.
\end{definition}

\end{frame}


\begin{frame}
\frametitle{A Matrix of Exponential Nonnegative Rank}

\pause
\begin{definition}
Let $M$ be a $2^n \times 2^n$ matrix, where each row (resp. column) is indexed by a $n$-bit string $a$ (resp. $b$), with
\[
M_{ab} \coloneqq {\left(1 - a^Tb\right)}^2
\]
\end{definition}

\pause
\begin{theorem}[De Wolf, '03]
Every $1$-monochromatic rectangle cover of $\suppmat(M)$ has size $2^{\Omega(n)}$.
\end{theorem}

\end{frame}

\begin{frame}
\frametitle{$\TSP(n)$ Polytope}

\pause
\begin{definition}[$\TSP$ Polytope]
Let $K_n = (V_n, E_n)$ be the complete graph with $n$ vertices. The \emph{Traveling Salesman Problem polytope} $\TSP(n)$ is defined as the convex hull of the characteristic vectors of all $F \subseteq E_n$ such that $F$ is a Hamiltonian cycle of $K_n$. In other words,
\[
\TSP(n) \coloneqq \conv\left(\left\{\chi^F \in \R^{E_n} \mid F \subseteq E_n \text{ is a Hamiltonian cycle of $K_n$}\right\}\right).
\]
\end{definition}

\pause
\begin{itemize}
\item Intuitively, the $\TSP(n)$ polytope is the polytope that has every Hamiltonian cycle of $K_n$ as a vertex.
\end{itemize}

\end{frame}

\begin{frame}
\frametitle{$\TSP(n)$ Polytope}

\pause
\begin{lemma}
For every $n$, there exists a positive integer $q = O(n^2)$ such that $\COR(n)$ is contained in a face of $\TSP(q)$.
\end{lemma}

\pause
\begin{theorem}
The extension complexity of $\TSP(n)$ is $\xc\left(\TSP(n)\right) = 2^{\Omega\left(\sqrt{n}\right)}$.
\end{theorem}

\end{frame}

\begin{frame}
\begin{center}
\LARGE{QUESTIONS ?}
\end{center}
\begin{figure}
        \centering
        \begin{minipage}{.5\textwidth}
            \centering
            \includegraphics[width=.8\linewidth]{questions.png}
        \end{minipage}%
        
\end{figure}
\end{frame}

\end{document}
